% Generated draft: requires mathematical and source review.
\documentclass[11pt]{article}
\usepackage[margin=25mm]{geometry}
\usepackage{fontspec}
\setmainfont{Latin Modern Roman}
\usepackage{amsmath,amssymb,mathtools,unicode-math}
\setmathfont{Latin Modern Math}
\usepackage{enumitem,needspace,xurl,hyperref}
\setlength{\parindent}{0pt}
\setlength{\parskip}{5pt}
\setlength{\emergencystretch}{3em}
\newcommand{\N}{\mathbb{N}}
\newcommand{\Z}{\mathbb{Z}}
\newcommand{\Q}{\mathbb{Q}}
\newcommand{\R}{\mathbb{R}}
\newcommand{\C}{\mathbb{C}}
\newcommand{\F}{\mathbb{F}}
\newcommand{\A}{\mathbb{A}}
\newcommand{\PP}{\mathbb{P}}
\newcommand{\E}{\mathbb{E}}
\newcommand{\eps}{\varepsilon}
\newcommand{\dd}{\,\mathrm{d}}
\begin{document}
\section*{UnsolvedMath ID: 6700059}
\textbf{Source problem code:} catalogue entry
\subsection{Definitions and mathematical statement}
For a convex polyhedron \(P\subset\R^n\), write \(Q_i\) for its facets and \(Q_{ij}=Q_i\cap Q_j\) for a codimension two face where facets meet. A combinatorial type is an isomorphism class of face incidence posets. The first task is to identify the types \(\mathcal P_{\mathrm{extr}}\) all of whose convex representatives are extremal, and to describe the individual extremal representatives of types \(\mathcal P_{\mathrm{nonextr}}\) which do not have that property. The antecedent definition of ``extremal'' is absent from this excerpt; it must not automatically be replaced by the following, newly introduced notion.

The source next defines \(P\) to be \emph{mean convexly extremal} if there is no region \(P'\subset\R^n\), diffeomorphic to \(P\) with corresponding faces \(Q'_i\), satisfying all of the following:
\begin{enumerate}
\item Every face has nonnegative mean curvature: \(H_{Q'_i}(p)\geq0\) at every point where this is defined.
\item For every incident pair \(i,j\) and every \(p\in Q'_{ij}=Q'_i\cap Q'_j\),
\[
 \theta'_{ij}(p)\leq\theta_{ij}(P),
\]
where \(\theta_{ij}(P)\) is the constant dihedral angle of \(P\), and \(\theta'_{ij}(p)\) is the corresponding angle between tangent facets of \(P'\).
\item For at least one pair \(i,j\) and at least one \(p_0\in Q'_{ij}\),
\[
 \theta'_{ij}(p_0)<\theta_{ij}(P).
\]
\end{enumerate}
The competitor is not stated to be convex or to have planar faces. No equality or inequality for face lengths, areas, or induced metrics is stipulated. The correspondence of faces is required to interpret the inequalities, but its precise diffeomorphism-with-corners regularity is not given. The mean curvature sign convention and the interior versus unoriented tangent-space angle convention are also missing; both must be fixed consistently before applying this definition. The two requests about ordinary extremality and the subsequent definition of mean convex extremality are retained as separate parts of the excerpt.
\subsection{Short English statement}
Classify combinatorial types whose every convex realization is extremal, and extremal realizations of the remaining types. The excerpt then defines mean convex extremality by forbidding mean convex competitors with all corresponding dihedral angles no larger and at least one strictly smaller.
\subsection{Sources}
\begin{itemize}
\item Local UnsolvedMath registry, permanent ID 6700059, source code catalogue entry; Misha Gromov, 2017; updated 2026-08-06.
\item Gromov, \emph{101 Questions, Problems and Conjectures around Scalar Curvature} (2017), printed tag [?63], occurrence 62, PDF page 60: \url{https://www.ihes.fr/~gromov/wp-content/uploads/2018/08/101-problemsOct1-2017.pdf}. Registry access date 2026-07-29; normalized PDF extraction; public author-hosted incomplete PDF; current status and reuse terms require review.
\item The summary associates the question with polyhedral scalar curvature and Alexandrov comparison but supplies no exact additional reference. Adjacent local ID 6700060 separately asks whether ordinary extremality implies mean convex extremality, confirming that the two terms should be distinguished but not defining the former.
\end{itemize}
\end{document}
